\documentclass[10pt,a4paper]{amsart}
\usepackage[margin=19mm]{geometry}
\usepackage{amsmath,amssymb,mathtools}
\usepackage[scr=rsfs,cal=euler]{mathalfa}
\usepackage{xfrac}
\usepackage[unicode,hidelinks,bookmarks=false]{hyperref}
\newcommand{\ie}{i.e.\ }
\newcommand{\cf}{cf.\ }
\newcommand{\resp}{respectively\ }
\newcommand{\Subsection}{\subsection}
\providecommand{\nobreakdash}{\nobreak\hbox{-}}
\setlength{\emergencystretch}{2em}
\multlinegap0pt
\usepackage{bm}
\usepackage{bbm}
\usepackage{dsfont}
\usepackage{xspace}
\usepackage{cancel}
\usepackage{mathtools}
\usepackage{amsthm}
\makeatletter
\@ifundefined{definition}{\newtheorem{definition}{Definition}}{}
\@ifundefined{lemma}{\newtheorem{lemma}[definition]{Lemma}}{}
\@ifundefined{proposition}{\newtheorem{proposition}[definition]{Proposition}}{}
\@ifundefined{theorem}{\newtheorem{theorem}[definition]{Theorem}}{}
\makeatother
\newcommand\B{\mathbb{B}}
\newcommand\HH{\mathcal{H}}
\newcommand\N{\mathbb{N}}
\newcommand\R{\mathbb{R}}
\newcommand\Z{\mathbb{Z}}
\newcommand\mcaF{\mathcal{F}}
\newcommand\mcaR{\mathcal{R}}
\title[Counterexamples to the $L^1$ and $L^{\infty}$ boundedness of]{Counterexamples to the $L^1$ and $L^{\infty}$ boundedness of the one-dimensional wave operators}
\author{Sisi Huang, Xiaohua Yao}
\date{Source: arXiv:2606.17898v1 (CC BY 4.0)}
\begin{document}
\maketitle

\noindent\emph{Source and licence.} Counterexamples to the $L^1$ and $L^{\infty}$ boundedness of the one-dimensional wave operators. Version arXiv:2606.17898v1 (CC BY 4.0), distributed under the Creative Commons Attribution 4.0 International licence, as recorded in the arXiv licence metadata for this exact version and shown on the abstract page https://arxiv.org/abs/2606.17898v1. Full rights notice: licenses/arxiv-2606.17898-CC-BY-4.0.txt.\par\smallskip

\noindent\textit{Editorial scope.} Counted unit: the Proposition whose source label is \texttt{good1} in the article (source0.tex lines 374--376, source label \texttt{good1}), delivered together with the article's own complete original proof of it (source0.tex lines 377--469, one \texttt{proof} environment of 93 lines). No mathematics is written, repaired or re-derived: statement and proof are the article's own text line for line. Required context retained uncounted: estim of gamma (lines 279--281); theorem of lp bounds of low energy (lines 293--323); kernels in low energy (lines 313--321); cancel lemma (lines 338--347); shur test (lines 369--371). Every definition, display and label of those lines is retained unchanged. Omitted from this package and not redistributed: 1--278, 282--292, 322--337, 348--368, 372--373, 470--1154. Nothing inside a retained excerpt is omitted or altered: every retained body is delivered exactly as the article's lines are written, so no omission receipt of any kind accompanies this package. Citations: the retained excerpts carry no citation command at all, so no bibliography entry of the article is transcribed and no reference is redistributed. Reference adaptation: none was needed and none was made; every reference inside the delivered unit resolves inside it, so no unresolved reference or citation is produced. Printed slips kept literal: no printed word, symbol, hypothesis or number of the counted statement or of its proof was altered. Layout and class: the article's own class is \texttt{amsart} with packages including color, amsmath, amssymb, euscript, mathrsfs, paralist, amsthm, ulem, bm, extarrows, verbatim, cite, hyperref, esint; the wrapper is the lane's pinned \texttt{amsart} editorial wrapper at 10pt on A4, which restates the theorem-like environments the retained text needs and adds the article's own macro definitions that the retained text needs (\texttt{\textbackslash B}, \texttt{\textbackslash HH}, \texttt{\textbackslash N}, \texttt{\textbackslash R}, \texttt{\textbackslash Z}, \texttt{\textbackslash mcaF}, \texttt{\textbackslash mcaR}), with extra wrapper packages (bm, bbm, dsfont, xspace, cancel, mathtools). The delivered numbering is the unit's own. Assets: no retained span contains a figure environment, a graphics-inclusion command, a PDF-page inclusion, a table or a TikZ picture; no image file is copied into this package, the package declares no asset dependency and no image file is redistributed with it. Licence: version arXiv:2606.17898v1 of this article is distributed under the Creative Commons Attribution 4.0 International licence, as recorded in the arXiv licence metadata for this exact version and shown on the abstract page https://arxiv.org/abs/2606.17898v1. A full rights notice is delivered at licenses/arxiv-2606.17898-CC-BY-4.0.txt. Characters: every retained excerpt is pure ASCII, so the delivered engine, XeLaTeX with its default Latin Modern text font, reports no missing character.
\begin{equation}\label{estim of gamma}
\|(\partial^{s}_{\lambda}\Gamma^{j}_{2})(\lambda)\|_{\B(L^2)}\leq C_s\lambda^{2-s},\quad 0<\lambda<2\lambda_0,\quad j=0,1.
\end{equation}

\begin{theorem}\label{theorem of lp bounds of low energy} 
Let $H=-\Delta+V$ with $|V(x)|\lesssim \left<x\right>^{-\beta}$ for some $\beta>2$. Then the following statements hold.    \begin{itemize}
    \item [{\rm(i)}] If $V$ is generic type and $\beta>7$, then 
\begin{equation}\label{deco of regular}
W^{L}_{-}=\sum\limits_{A\in\Omega_0 }K_{A}+K_{{P_1}}
\end{equation}
and $\{K_A:A\in\Omega_0\}\subseteq\B(L^p)$ for all $1\leq p\leq\infty$, while 
$$ K_{{P_1}}\in \B(L^p)\cap\B(\HH^1,L^1)\cap\B(L^{\infty},{\rm BMO})\ for\ all\ 1<p<\infty.$$
Therefore, $$W^L_{-}\in\B(L^p)\cap\B(\HH^1,L^1)\cap\B(L^{\infty},{\rm BMO})\ for\ all\ 1<p<\infty.$$ 
\item [{\rm(ii)}] If $V$ is exceptional type and $\beta>11$, then 
\begin{equation}\label{deco of resonance}
W^{L}_{-}=\sum\limits_{A\in\Omega_1}K_A+K_{-1}+\sum\limits_{j=1}^{2}(K_{0j}+K_{{P_{j}}})   
\end{equation}
and $\{K_A:A\in\Omega_1\}\subseteq\B(L^p)$ for all $1\leq p\leq\infty$, while  
\begin{align*}
K_{-1}&\in\B(L^p)\cap\B(\HH^1,L^1)\ for\ all\ 1<p<\infty,\\
\{K_{01},K_{02},K_{{P_1}},K_{{P_2}}\}&\subseteq \B(L^p)\cap\B(\HH^1,L^1)\cap\B(L^{\infty},{\rm BMO})\ for\ all\ 1<p<\infty.
\end{align*}
Therefore,  $$W^L_{-}\in\B(L^p)\cap\B(\HH^1,L^1)\ for\ all\ 1<p<\infty.$$ 
Here $\Omega_j=\{QA^j_0Q,\lambda QA^j_{11},\lambda A^j_{12}Q, \Gamma^j_2(\lambda)\}$ and
\begin{align}\label{kernels in low energy}
\begin{split}
K_A(x,y)&=\int_{0}^{+\infty}\lambda\chi(\lambda)\big[R^+_0(\lambda^{2})vAv(R^+_0-R^-_0)(\lambda^{2})\big](x,y)d\lambda,\quad A\in\Omega_j,\quad j=0,1,\\
K_{-1}(x,y)&=-ic\int_{0}^{+\infty}\chi(\lambda)\big[R^+_0(\lambda^{2})vS_1v(R^+_0-R^-_0)(\lambda^{2})\big](x,y)d\lambda,\\
K_{01}(x,y)&=2\alpha c\int_{0}^{+\infty}\lambda\chi(\lambda)\big[R^+_0(\lambda^{2})vS_1T_0v(R^+_0-R^-_0)(\lambda^{2})\big](x,y)d\lambda,\\
K_{02}(x,y)&=2\alpha c\int_{0}^{+\infty}\lambda\chi(\lambda)\big[R^+_0(\lambda^{2})vT_0S_1v(R^+_0-R^-_0)(\lambda^{2})\big](x,y)d\lambda,\\
K_{P_{j}}(x,y)&=\int_{0}^{+\infty}\lambda^2\chi(\lambda)\big[R^+_0(\lambda^{2})vP_jv(R^+_0-R^-_0)(\lambda^{2})\big](x,y)d\lambda,\quad j=1,2.
\end{split}
\end{align}
\end{itemize}
\end{theorem}

\begin{lemma}\label{cancel lemma}
Let $\lambda>0$. Then for any $T=Q,S_1$, we have
\begin{equation*}
 [R^{\pm}_0(\lambda^2)vTf](x)=\int_{\R}\mcaR^{\pm}(\lambda,x,y)(Tf)(y)dy,\quad\ [TvR^{\pm}_0(\lambda^2)f](x)=T(\mcaF^{\pm}(\lambda,x)),
\end{equation*}
where 
\begin{equation*}
\mcaR^{\pm}(\lambda,x,y)=\frac{1}{2}yv(y)\int_{0}^{1}{\rm sgn}(x-\theta y)e^{\pm i \lambda|x-\theta y|}d\theta,\quad \mcaF^{\pm}(\lambda,x)=\int_{\R}\mcaR^{\pm}(\lambda,y,x)f(y)dy.
\end{equation*}
\end{lemma}

\begin{lemma}\label{shur test}
If the kernel $K(x,y)$ satisfies $$\sup_{x\in\R}\int|K(x,y)|dy+\sup_{y\in\R}\int|K(x,y)|dx<\infty,$$ then $K\in \B(L^{p}(\R))$ for all $1\leq p\leq \infty.$
\end{lemma}

\begin{proposition}\label{good1}
Under the assumption of Theorem {\rm\ref{theorem of lp bounds of low energy}}, let $\Omega_j$ and $K_A$ be as in \eqref{kernels in low energy}. Then for any $j=0,1$ and $A\in\Omega_j$, we have $K_A\in\B(L^p)$ for all $\ 1\leq p\leq\infty$.
\end{proposition}

\begin{proof} Recall that $\Omega_j=\{QA^j_0Q,\lambda QA^j_{11},\lambda A^j_{12}Q,  \Gamma^j_2(\lambda)\}$, it suffices to consider the case $A\in\Omega_0$. Moreover, we only handle $K_A$ for $A=QA^0_0Q$ and $A=\Gamma^0_2(\lambda)$, since the remained cases are similar.
\vskip0.2cm
\underline{(1) For $A=QA^0_0Q$}, denote
\begin{equation*}
K_0(x,y)=\int_{0}^{+\infty}\lambda\chi(\lambda)\big[R^+_0(\lambda^{2})vQA^0_0Qv(R^+_0-R^-_0)(\lambda^{2})\big](x,y)d\lambda.
\end{equation*}
By Lemma \ref{cancel lemma}, we have 
\begin{equation*}
    K_0(x,y)=\frac{1}{4}\sum\limits_{\pm}K^{\pm}_0(x,y),
\end{equation*}
where $\Theta=(u_1,u_2,\theta_1,\theta_2),\ X_1=x-\theta_1u_1,\ Y_2=y-\theta_2u_2$ and 
\begin{align*}
K^{\pm}_0(x,y)&=\pm\int_{0}^{+\infty}e^{i\lambda(|x|\pm|y|)}\lambda\chi(\lambda)k^{\pm}_0(\lambda,x,y)d\lambda,\\
k^{\pm}_0(\lambda,x,y)&=\int_{\R^2\times[0,1]^2}({\rm sgn}X_1)({\rm sgn}Y_2)e^{i\lambda(|X_1|-|x|\pm (|Y_2|-|y|))}u_1u_2(vQA^0_0Qv)(u_1,u_2)d\Theta.
\end{align*}
We shall prove \begin{equation}\label{esti of Kpm0}
    |K^{\pm}_{0}(x,y)|\lesssim\left<|x|-|y|\right>^{-2},\quad \forall\  x,y\in \R.
\end{equation}
Together with Lemma \ref{shur test} this yields $K_0\in \B(L^p)$ for all $1\leq p\leq\infty$. First, we note that $K^{\pm}_0(x,y)$ is uniformly bounded since $\chi(\lambda)\in C^{\infty}_0(\R)$ and 
\begin{equation*}
|k^{\pm}_0(\lambda,x,y)|\lesssim \int_{\R}\left<u_1\right>|v(u_1)|\cdot |QA^0_0Q|(\left<\cdot\right>|v(\cdot)|)(u_1)du_1 \lesssim\|\left<u_1\right>^2V(u_1)\|_{L^1} <\infty. 
\end{equation*}
Thus, if $||x|-|y||\leq1$, then $|K^{\pm}_{0}(x,y)|\lesssim1\lesssim\left<|x|-|y|\right>^{-2}.$ Now suppose $||x|-|y||\geq1$ and set $$\Phi^{\pm}(x,y,\Theta)=i(|X_1|-|x|\pm (|Y_2|-|y|)).$$
By the triangle inequality, $\left|\Phi^{\pm}(x,y,\Theta)\right|\leq |u_1|+|u_2|$. Hence, for any $\ell=0,1,2$, 
\begin{equation}\label{estimate of k_0}
|(\partial^{\ell}_{\lambda}k^{\pm}_0)(\lambda,x,y)|\lesssim
\int_{\R^2}^{}\left<u_1\right>^{\ell+1}|(vQA^0_0Qv)(u_1,u_2)|\left<u_2\right>^{\ell+1}du_1du_2\lesssim \|\left<u_1\right>^{2\ell+2}V(u_1)\|_{L^1} <\infty.
\end{equation} 
Let $L^{\pm}(\lambda,x,y)=\lambda \chi(\lambda)k^{\pm}_0(\lambda,x,y)$. Integrating by parts twice gives
\begin{align}\label{inte parts twice of Kpm0}
\begin{split}
K^{\pm}_{0}(x,y)&=\pm\left[\frac{e^{i\lambda(|x|\pm |y|)}}{i(|x|\pm |y|)}L^{\pm}(\lambda,x,y)\Big|^{+\infty}_{\lambda=0}-\int_{0}^{+\infty}\frac{e^{i\lambda(|x|\pm |y|)}}{i(|x|\pm |y|)}(\partial_{\lambda}L^{\pm})(\lambda,x,y)\right]\\
&=\frac{\pm}{(|x|\pm |y|)^2}\left(e^{i\lambda(|x|\pm |y|)}(\partial_{\lambda}L^{\pm})(\lambda,x,y)\Big|^{+\infty}_{\lambda=0}-\int_{0}^{+\infty}e^{i\lambda(|x|\pm |y|)}(\partial^2_{\lambda}L^{\pm})(\lambda,x,y
)d\lambda\right)\\
&=O(\left<|x|-|y|\right>^{-2}),
\end{split}
\end{align}
where the second equality uses $L^{\pm}(\lambda,x,y)\big|_{\lambda=0,+\infty}=0$, and the last step follows from $\chi(\lambda)\in C^{\infty}_0(\R)$ and \eqref{estimate of k_0}.
This establishes \eqref{esti of Kpm0}.
\vskip0.3cm
\underline{(2) For $A=\Gamma^0_2(\lambda)$}, let 
\begin{equation*}
K_r(x,y)=\int_{0}^{+\infty}\lambda\chi(\lambda)\big[R^+_0(\lambda^{2})v\Gamma^0_2(\lambda)v(R^+_0-R^-_0)(\lambda^{2})\big](x,y)d\lambda.
\end{equation*}
We can write
\begin{equation*}
    K_r(x,y)=-\frac{1}{4}\sum\limits_{\pm}K^{\pm}_r(x,y),
\end{equation*}
where $X_1=x-u_1,\ Y_2=y-u_2,\ {\Gamma}(\lambda)=\lambda^{-2}\Gamma^0_2(\lambda)$ and 
\begin{align*}
K^{\pm}_r(x,y)&=\int_{0}^{+\infty}e^{i\lambda(|x|\pm|y|)}\lambda\chi(\lambda)k^{\pm}_r(\lambda,x,y)d\lambda,\\
k^{\pm}_r(\lambda,x,y)&=\int_{\R^2}e^{i\lambda(|X_1|-|x|\pm (|Y_2|-|y|))}(v\Gamma(\lambda)v)(u_1,u_2)du_1du_2.
\end{align*}
\indent Compared with $K_0$, the additional complexity comes from $k^{\pm}_r$. To handle it we introduce a homogeneous dyadic partition of unity $\{{\varphi_{N}}\}_{N\in \Z}$ on $(0,+\infty)$: choose $\varphi\in C^{\infty}_{0}(\R^{+})$ with $0\leq\varphi\leq1$, supp$~\varphi\subset[\frac{1}{4},1],$ set $\ \varphi_{N}(\lambda)=\varphi(2^{-N}\lambda)$ so that supp$~\varphi_{N}\subset[2^{N-2},2^{N}]$, and $$ \sum_{N\in\Z}{}\varphi_{N}(\lambda)=1,  \quad\lambda>0.$$
Let $N_0=[\log_22\lambda_0]$. Combining the supports of $\chi(\lambda)$ and $\varphi_N(\lambda)$, we decompose
\begin{equation*}
\chi(\lambda)=\sum_{N=-\infty}^{N_0}{\chi}_{N}(\lambda),\quad {\chi}_{N}(\lambda)=\chi(\lambda)\varphi_{N}(\lambda),\quad\lambda>0.
\end{equation*}
Consequently,
\begin{equation*}
K^{\pm}_{r}(x,y)=\sum_{N=-\infty}^{N_0}\int_{0}^{+\infty}e^{i\lambda(|x|\pm |y|)}\lambda{\chi}_{N}(\lambda)k^{\pm}_r(\lambda,x,y)d\lambda:=\sum_{N=-\infty}^{N_0}K^{\pm}_{r,N}(x,y).
\end{equation*}
We claim that for every  $N\leq N_0$,
\begin{equation}\label{estim of krlN}
|K^{\pm}_{r,N}(x,y)|\lesssim \min\{{2^{2N},\left<|x|-|y|\right>^{-2}}\},\quad \forall\ x,y\in\R.
\end{equation}
Once this is proved, it implies for any $\theta\in[0,1]$, 
 $$|K^{\pm}_{r,N}(x,y)|\lesssim2^{2N(1-\theta)}\left<|x|-|y|\right>^{-2\theta},\quad \forall\ x,y\in\R.$$
Choosing $\theta=\frac{3}{4}$ gives
$$|K^{\pm}_{r}(x,y)|\lesssim\sum_{N=-\infty}^{N_0}2^{\frac{N}{2}}\left<|x|-|y|\right>^{-\frac{3}{2}}\lesssim\left<|x|-|y|\right>^{-\frac{3}{2}},\quad \forall\ x,y\in\R,$$
and Lemma \ref{shur test} yields $K_r\in\B(L^p)$ for all $1\leq p\leq\infty$. To derive \eqref{estim of krlN}, we first recall from \eqref{estim of gamma} that
$$\|\Gamma^{(j)}(\lambda)\|_{\B(L^2)}=O(\lambda^{-j}),\quad 0<\lambda\leq2\lambda_0.$$
This gives, for any $\ell=0,1,2$,
\begin{equation}\label{estim of krN}
\sup_{x,y\in\R}\left|(\partial^{\ell}_{\lambda}k^{\pm}_r)(\lambda,x,y)\right|\lesssim\lambda^{-\ell}\big\|\left<\cdot\right>^{2\ell}V(\cdot)\big\|_{L^1(\R)},\quad  0<\lambda\leq2\lambda_0.
\end{equation}
Now fix $ N\leq N_0$. Because supp$\ \chi_N\subseteq [0,2\lambda_0]\cap[2^{N-2},2^N]$, 
\begin{align}\label{klbn}
|K^{\pm}_{r,N}(x,y)|=\left|\int_{0}^{+\infty}e^{i\lambda(|x|\pm |y|)}\lambda{\chi}_{N}(\lambda)k^{\pm}_r(\lambda,x,y)d\lambda\right|\lesssim\int_{2^{N-2}}^{2^{N}}\lambda d\lambda\lesssim2^{2N},\quad \forall \ x,y\in\R.
\end{align}
Hence, if $||x|-|y||\leq1$, then
$$|K^{\pm}_{r,N}(x,y)|\lesssim2^{2N}\lesssim2^{2N_0} \left<|x|-|y|\right>^{-2}\lesssim \left<|x|-|y|\right>^{-2},\quad \forall\ N\leq N_0 ,$$
If $||x|-|y||\geq1$, set $L^{\pm}_{N}(\lambda,x,y)=\lambda{\chi}_{N}(\lambda)k^{\pm}_r(\lambda,x,y)$.
Apply integration by parts twice to $K^{\pm}_{r,N}(x,y)$ and using ${\chi}^{(\ell)}_N(\lambda)\big|_{\lambda=0,+\infty}=0$ together with \eqref{estim of krN}, we have
\begin{align*}
\big|K^{\pm}_{r,N}(x,y)\big|=\frac{1}{(|x|\pm|y|)^2}\left|\int_{0}^{+\infty}e^{i\lambda(|x|\pm |y|)}(\partial^2_{\lambda}L^{\pm}_{N})(\lambda,x,y)d\lambda\right|\lesssim\frac{1}{(|x|\pm |y|)^2}\int_{2^{N-2}}^{2^N}\big|(\partial^2_{\lambda}L^{\pm}_{N})(\lambda,x,y)\big|d\lambda.
\end{align*}
Since
$\big|{\chi}^{(s)}_{N}(\lambda)\big|\lesssim2^{-Ns}$ for any $s\in\N$,
this bound together with \eqref{estim of krN} implies 
$$\big|(\partial^2_{\lambda}L^{\pm}_{N})(\lambda,x,y)\big|\lesssim 2^{-N},\quad \lambda\in[2^{N-2},2^N]$$
and therefore $\big|K^{\pm}_{r,N}(x,y)\big|\lesssim \left<|x|-|y|\right>^{-2}.$ The claim \eqref{estim of krlN} follows.
\end{proof}

\end{document}
